\clearpage
\fancyhead[L]{\small GAUSS AN BODE.}
\hypertarget{native-289-lower}{}
\begin{flushright}
Braunschweig, den 23. August 1807.
\end{flushright}

\par\indent Verzeihen Sie bitte, dass ich Ihnen seit der Mitteilung meiner ersten Elemente der Vesta keine weiteren Nachrichten über meine späteren Untersuchungen zu diesem Gegenstand geschickt habe.
\SourceLeaf{290}%
\par\indent Eine Reise nach Bremen in diesem Sommer und eine Krankheit, die mich danach traf, haben mich sowohl in meinen Arbeiten insgesamt als auch in meiner Korrespondenz etwas zurückgeworfen. Jetzt beeile ich mich umso mehr, das Versäumte nachzuholen, damit Sie meine Mitteilungen noch nach Belieben in Ihrem neuen Jahrbuch verwenden können.

\par\indent Meine zweiten Elemente der \VestaBodyTwoNinety{} und ihre Übereinstimmung mit sehr vielen Beobachtungen werden Ihnen bereits aus der \href{https://zs.thulb.uni-jena.de/rsc/viewer/jportal_derivate_00237836/Monatliche_Correspondenz_130168688_15_1807_0605.tif}{M.\,C.} bekannt sein.

\par\indent Für die Berechnung der \emph{dritten} Elemente habe ich außer den früheren Beobachtungen auch die Pariser und Mailänder Meridianbeobachtungen und die letzten Bremer und Lilienthaler Beobachtungen --- sie reichen bis zum 11. Juli --- verwendet. Hier sind diese dritten Elemente.

\begin{center}\begingroup\fontsize{9}{11}\selectfont
\begin{tabular}{ll}
\shortstack[l]{Epoche der Länge, 31. März 1807\\$0^{\text{u}}$ in Bremen} .\ .\ . & $192^\circ\,23'\,30''1$ \\
Aphel .\ .\ .\ .\ .\ .\ .\ .\ .\ .\ .\ .\ .\ .\ .\ . & $69\;50\;31.9$ \\
Knoten .\ .\ .\ .\ .\ .\ .\ .\ .\ .\ .\ .\ .\ .\ .\ .\ .\ . & $103\;18\;28.0$
\end{tabular}\endgroup\end{center}

\noindent Beide gelten für die oben genannte Epoche und werden als relativ zu den Fixsternen unbeweglich vorausgesetzt.

\begin{center}\begingroup\fontsize{9}{11}\selectfont
\begin{tabular}{ll}
Neigung der Bahn .\ .\ .\ .\ .\ .\ .\ .\ .\ .\ .\ .\ . & $7\;8\;10.7$ \\
Exzentrizität .\ .\ .\ .\ .\ .\ .\ .\ .\ .\ .\ .\ .\ .\ . & $0.0855050$ \\
Logarithmus der Halbachse .\ .\ .\ .\ .\ .\ .\ .\ .\ . & $0.3720160$ \\
Mittlere tägliche tropische Bewegung .\ .\ .\ .\ .\ . & $981''7087$
\end{tabular}\endgroup\end{center}

\par\indent Ich füge hier auch die nach diesen Elementen berechnete Ephemeride für das nächste Jahr bei.

[Rektaszension und Deklination sind auf Minuten genau für Mitternacht in Seeberg in Abständen von vier Tagen vom 2. Mai 1808 bis zum 18. März 1809 angegeben.]

